\chapter{Lazy initialisation}
\label{ch:lazy}

\begin{aside}
Startup is a precious window. Anything you can defer makes
the app feel snappier. Lazy initialisation pushes work to
``when needed'', not ``at start''.
\end{aside}

\section{What to defer}

\begin{itemize}[leftmargin=*]
  \item Loading large datasets.
  \item Connecting to network services.
  \item Compiling regexes.
  \item Reading large files.
  \item Initialising plug-ins.
\end{itemize}

\section{Lazy signals}

A signal whose value is computed on first read:

\begin{lstlisting}
auto big_data = lazy_signal([&] {
    return load_from_disk();
});
\end{lstlisting}

The load runs when something reads the signal, not at
declaration.

\section{Lazy widgets}

A tab that isn't shown doesn't need to be rendered. Defer
its construction:

\begin{lstlisting}
if (ctx.get(active_tab) == 2) {
    return analytics_tab();
}
\end{lstlisting}

The tab's widgets and signals only exist when active.

\section{Cold start cost}

Lazy means first use is slower. For predictable cold-start
cost, prefetch in the background after startup completes:

\begin{lstlisting}
ctx.on_idle([&] {
    prefetch(big_data);
});
\end{lstlisting}

The user sees a fast startup; the background does the
loading.

\section{Caching}

A lazy value, once computed, caches. Subsequent reads are
fast. Invalidate via explicit reload.

\section{Recap}

\begin{recap}
\begin{itemize}[leftmargin=*]
  \item Defer expensive work past startup.
  \item Lazy signals; lazy widgets.
  \item Idle prefetch to smooth first-use latency.
  \item Cache the result of a lazy compute.
\end{itemize}
\end{recap}

\section*{Workshop}

\begin{enumerate}[leftmargin=*]
  \item Identify your app's startup hot path; defer
    something.
  \item Add an idle prefetch for a heavy resource.
  \item Measure startup time before and after.
\end{enumerate}
